% Nouveau EGA V (cote 25): archive cote 25, site page 86 (PDF page 87).
% typescript. Draft compared: none.
% AI-transcribed by Claude Opus 5.5 (claude-opus-5-5, adaptive thinking, effort high), Message Batches API, batch ega25b (msgbatch_01A7n5YuXubUPNK1zsepi8RG), 2026-11-13; checked against both supplied scan parts; see CHECK_DECISIONS.tsv.
\scanpage{25/86}
\mspage{\hand{\unsure{80}}}

p$_{\var{\deleted{\unsure{C}}}{}{deletion}{ega25_25_086-v0001}}$ précédent, \var{\deleted{qu'il sera commode de noter également de la même façon}}{}{deletion}{ega25_25_086-v0002}
\var{\deleted{p$_C$ , pour indiquer sa dépendance de C.}}{}{deletion}{ega25_25_086-v0003} De cette façon, les propriétés de
la projection conique  restreinte  sont ramenées \var{\deleted{dans une large mesure}}{}{deletion}{ega25_25_086-v0004} à
celles de la projection conique étendue, qui a été  étudiée systématique-%
ment par ailleurs, ou est sensée l'avoir été (cf n\textsuperscript{o} 10 et n\textsuperscript{o} 12). La
question\var{\xout{s}}{}{typewriter strike-out}{ega25_25_086-v0005} principale\var{\xout{s}}{}{typewriter strike-out}{ega25_25_086-v0006} qui se pose est, lorsque  S=Spec(k), quelles sont
les propriétés de la p\unsure{ro|oo}jection conique  de X  quand on prend C générique
dans Grass$^{m+1}$(\unsure{E|P}) (ce qui exige  qu'on fasse un changement de base
k$\hand{\rightarrow}$k(\hand{$\unsure{\eta|y}$}) , i.e.  C est \var{\inserted[typed]{alors}}{alors}{insertion}{ega25_25_086-v0007} en fait une sous-variété linéaire de X$_{k(\hand{\unsure{\eta|y}})}$ ); des
\unsure{ar|aa}guments  standart, déjà répétés à satiété, permettront d'en conclure
les propriétés analogues pour les projections coniques correspondant à \var{\xout{de}}{}{typewriter strike-out}{ega25_25_086-v0008}
des points de Grass$^{m+1}$(P)  appartenant à un ouvert non vide de ladite
grassmanienne, et enfih, lorsque k est infini, on conclut à l'existence
d'une(en fait une infinité) \var{\inserted[typed]{de}}{de}{insertion}{ega25_25_086-v0009} \var{\xout{de}}{}{typewriter strike-out}{ega25_25_086-v0010} C  \uline{définie\var{\xout{s}}{}{typewriter strike-out}{ega25_25_086-v0011}} \uline{sur} k , I.e.  une sous-variété
linéaire de P lui-même (sans changement de corps de base), donnant lieu
à une projection conique ayant les propriétés en question. Il y aurait
\var{\margin{|}}{\placedmargin{|}}{marginal note placed}{ega25_25_086-v0012} lieu de bloquer ce genre d'explications générales avec celles de la même
eau données dans les n\textsuperscript{os} 4 et 7, et que nous avons d'ailleurs déjà
utilisées plus ou moins implicitement,  par exemple au n\textsuperscript{o}13 \var{\xout{\unsure{(dém. de}}}{}{typewriter strike-out}{ega25_25_086-v0013}
\var{\xout{13.3}}{}{typewriter strike-out}{ega25_25_086-v0014} (cf 13.4. c). Il y a lieu également, à ce propos, d'examiner les
propriétés relatives à un faisceau F sur X,  en prenant son image inverse
$F^{(m)}_{\hand{Q}(C)}$  sur  $\underline{X}^{(m)}_{\hand{Q}(C)}$ . Il faut  d'ailleurs, dans la situation précise
envisagée ici, d es notations plus simples, je propose  $\overset{\hand{\sim}}{X}(C)$ et $\overset{\hand{\sim}}{F}(\unsure{C|O})$ ,
et simplement $\overset{\hand{\sim}}{X}$ et $\overset{\hand{\sim}}{F}$  quand aucune confusion n\unsure{'}en résulte  (attention, le
F \var{\deleted{précédent}}{}{deletion}{ega25_25_086-v0015} \var{\inserted{\unsure{ici}}}{\unsure{ici}}{insertion}{ega25_25_086-v0016} n'est pas  celui du début du n\textsuperscript{o}).

Grosso modo, et lorsque  disons on suppose  f une immersion, les
propriétés de la projection conique générique sont assez différentes,

\mspage{\hand{86}}
